\section{Introduction}
\label{sec:introduction}

EDA scripting is challenging because tool APIs are tool-specific, sparsely documented, and exposed through SWIG bindings whose type rules differ from native Python. ZhuLong~\cite{liu2026zhulong} is the first work that established an execution-grounded LLM agent for EDA scripting that combines API retrieval, documentation inspection, and sandbox execution via unified MCP tools, augmented by an offline API self-exploration mechanism that infers undocumented API behaviors through counterfactual experimentation. On the EDA-Eval-PyAether benchmark, the system achieved 78.5\% pass@1, substantially outperforming a pure LLM baseline of 23.6\%. However, a critical limitation remained: the agent solved each task independently and did not learn from past execution traces, so its performance on new tasks did not improve over time.
% EDA脚本编写面临挑战，因为工具API具有工具特异性、文档稀疏，且通过SWIG绑定暴露，其类型规则与原生Python不同。ZhuLong~\cite{liu2026zhulong}首次建立了一个基于执行反馈的EDA脚本LLM智能体，通过统一的MCP工具组合API检索、文档查阅和沙箱执行，并辅以离线API自探索机制，通过反事实实验推断未文档化的API行为。在EDA-Eval-PyAether基准上，该系统达到了78.5\%的pass@1，显著优于23.6\%的纯LLM基线。然而，一个关键局限仍然存在：该智能体独立解决每个任务，并未从过去的执行轨迹中学习，因此在新任务上的性能不会随历史经验积累而提升。

Building on this task-level substrate, we introduce \textbf{self-evolution}: an offline system that distills execution traces into reusable knowledge and re-injects it into future tasks, with reliability gates that prevent unverified experience from misleading the agent. To provide reliable signals for self-evolution, we introduce a precision-optimized static checker for deterministic API misuse. It raises no error flags on the 158 benchmark ground-truth solutions, a sanity check limited to this reference-code set. Applying self-evolution to EDA scripting introduces two challenges. First, feedback is very \textbf{sparse}: the agent receives only a binary pass/fail from the hidden verifier. This is different from methods like SkillOpt~\cite{yang2026skillopt}, which assume dense per-step rewards. Without careful verification, distilled experience can mislead. Second, instance-disjoint evaluation risks cross-split answer copying: a held-out task may copy a training sibling's verbatim solution rather than transferring the underlying method, a problem recently observed in EvolveMem~\cite{evolvemem2026}.
% 在这一任务级基础之上，我们引入\textbf{自进化}：一个离线系统将执行轨迹蒸馏为可复用知识并重新注入未来任务，通过可靠性门控防止未经验证的经验误导智能体。为给自进化提供可靠信号，我们引入了针对确定性API误用的精度优化静态检查器；在基准任务的ground-truth解上未观察到误报，此健全性检查仅限于该参考代码集。将自进化应用于EDA脚本编写带来了两个挑战。第一，反馈是\textbf{稀疏}的：智能体仅从隐藏验证器收到二元的通过/失败信号——与SkillOpt~\cite{yang2026skillopt}不同，后者假设密集的逐步奖励。缺乏审慎验证时，蒸馏经验可能误导。第二，实例不相交评估存在跨划分答案复制风险：留出任务可能直接复制训练集同类任务的逐字解决方案，而非迁移其方法，这一问题近期也在EvolveMem~\cite{evolvemem2026}中被观察到。

We present \textbf{ZhuLong-SE (Self-Evolution)}, which extends the ZhuLong task-level substrate with enhanced API retrieval and documentation responses, improved sandbox isolation, and precision-optimized static checking, and then adds a self-evolution layer with reliability gating. To handle sparse feedback, we introduce ground-truth anchoring, scope-boundary reflection, disconfirmation-based filtering, and a held-out validation loop that rejects degrading rounds. Beyond post-hoc evaluation gating, we design structural anti-leakage retrieval via names-only injection and deterministic docstring generalization to reduce the risk of cross-split answer copying. These mechanisms are designed to provide methodological signals without injecting full call sequences, while an evolution round is retained only when held-out validation does not degrade beyond the acceptance threshold. On the held-out test set, this design achieves a higher mean pass@1 than the human-crafted-experience baseline.
% 我们提出了\textbf{ZhuLong-SE（Self-Evolution）}，在ZhuLong任务级基础之上增强API检索与文档响应、改进沙箱隔离并加入精度优化静态检查，随后增加带可靠性门控的自进化层。为处理稀疏反馈，我们引入真值锚定、范围边界反思、基于证伪的过滤，以及拒绝退化轮次的留出验证循环。在事后评估门控之外，我们还通过仅名称注入和确定性文档字符串泛化设计了结构化反泄漏检索，以降低跨划分答案复制的风险。这些机制旨在提供方法论信号而不注入完整调用序列，同时仅在留出验证未超过退化阈值时保留进化轮次。在留出测试集上，该设计取得了高于人工经验基线的平均pass@1。

Our contributions are as follows.
% 我们的贡献如下。

\begin{itemize}[leftmargin=*]
\item \textbf{Self-evolution as a design methodology for EDA automation.} 
We establish self-evolution---distilling execution traces into reusable knowledge and re-injecting them into future tasks---as a systematic methodology for system-level EDA scripting. To make it deployable under EDA's sparse feedback and opaque commercial APIs, we introduce three governance rules: ground-truth anchoring, on-demand pitfall retrieval, and validation-gated rollback.
% 自进化作为EDA自动化的一种设计方法论。
% 我们将自进化——将执行轨迹蒸馏为可复用知识并重新注入到未来任务中——确立为一种系统级的EDA脚本编写方法论。为使其在EDA稀疏反馈和不透明商业API条件下可部署，我们引入了三条治理规则：真值锚定、按需陷阱检索、以及验证门控回滚。

\item \textbf{A rigorous evaluation protocol against cross-split leakage.} 
We expose cross-split answer copying as a threat to trustworthy self-evolution evaluation, and propose a retrieval-time defense (names-only injection + docstring generalization) that removes full call sequences and task-specific details while preserving methodological transfer---establishing a protocol for future work on self-evolving design agents.
% 一个对抗跨划分泄漏的严格评估协议。
% 我们揭示了跨划分答案复制对可信自进化评估构成的威胁，并提出了一种检索时防御机制（仅名称注入+文档字符串泛化），在保留方法论迁移信号的同时移除完整调用序列和任务特定细节——为未来自进化设计智能体的研究建立了一个评估协议。

\item \textbf{Counterintuitive design insights from controlled ablation.} 
Controlled ablations on the 158-task EDA benchmark yield three observations: ground-truth anchoring prevents 12 of 24 failed-task traces from producing falsely verified entries; pitfall pre-injection reduces validation pass@1 by 8.0\,pp in our ablation; and names-only retrieval improves over full-sequence injection by 4.7\,pp. On the 40-task held-out test set, self-evolution from the Bare configuration reaches 80.0\% mean pass@1, compared with 77.5\% for the human-crafted-experience baseline, grounding these observations for future system-level design automation.
% 来自控制变量消融实验的反直觉设计洞察。
% 在158任务EDA基准上的受控消融得到三项观察：真值锚定阻止24条失败任务轨迹中的12条产生错误验证条目；在当前消融中，陷阱预注入使验证pass@1下降8.0个百分点；仅名称检索相比完整序列注入提升4.7个百分点。在40任务留出测试集上，从Bare配置开始的自进化达到80.0\%平均pass@1，而人工经验基线为77.5\%，为未来系统级设计自动化提供了实证依据。
\end{itemize}
